\section{Rectificador de onda completa}

\subsection{Instrumentos y componentes utilizados}

\begin{itemize}
    \item Fuente de tensión alterna de 12 $V_{RMS}$, $V_S$, y 50 Hz, $f_S$.
    \item Osciloscopio.
    \item Multímetro digital.
    \item Puente rectificador compuesto por cuatro diodos, $D_1$, $D_2$, $D_3$ y $D_4$.
    \item Fusible de protección conectado en serie con la entrada, $F$.
    \item Resistencia limitadora de corriente de valor nominal $R_{lím} = 4{,}7\,\Omega$.
    \item Resistencia de carga de valor nominal $R_L = 100\,\Omega$.
    \item Tres capacitores electrolíticos: $C_1 = 100\,\mu$F, $C_2 = 220\,\mu$F, $C_3 = 2200\,\mu$F.
    \item Cables para el conexionado.
\end{itemize}


\subsection{Objetivos}

El objetivo de esta experiencia fue estudiar experimentalmente el funcionamiento de un circuito rectificador de onda completa y comparar su desempeño con el del rectificador de media onda analizado previamente. En particular, se buscó observar la señal de salida obtenida a partir de una entrada sinusoidal, determinar sus características mediante el osciloscopio y el multímetro, y estudiar la modificación de la tensión de salida al incorporar capacitores de filtrado de diferentes capacidades. Asimismo, se calculó analíticamente el factor de rizado para la señal sin capacitor y se lo comparó con el obtenido para el rectificador de media onda, con el fin de evaluar la eficiencia energética de ambas configuraciones.


\subsection{Graficos}

\begin{figure}[H]
    \centering
    \begin{tikzpicture}
        \begin{axis}[
            width=\textwidth,
            height=6cm,
            xlabel={Tiempo [ms]},
            ylabel={Tensi\'on [V]},
            title={Rectificador de onda completa --- $C_1 = 100\,\mu$F},
            grid=both,
            grid style={line width=0.3pt, draw=gray!30},
            xmin=0, xmax=45,
            ymin=5, ymax=18,
            xtick distance=10,
            ytick distance=5,
            tick label style={font=\small},
            label style={font=\small},
            title style={font=\normalsize\bfseries},
        ]
        \addplot[color=magenta!70!black, line width=0.8pt]
            table[col sep=comma, x=tiempo, y=voltaje]{CSV/onda completa/oc_C1_100uF.csv};
        \end{axis}
    \end{tikzpicture}
    \caption{Tensi\'on de salida del rectificador de onda completa con $C_1 = 100\,\mu$F.}
    \label{fig:oc_c1_100uF}
\end{figure}

\begin{figure}[H]
    \centering
    \begin{tikzpicture}
        \begin{axis}[
            width=\textwidth,
            height=6cm,
            xlabel={Tiempo [ms]},
            ylabel={Tensi\'on [V]},
            title={Rectificador de onda completa --- $C_2 = 220\,\mu$F},
            grid=both,
            grid style={line width=0.3pt, draw=gray!30},
            xmin=0, xmax=45,
            ymin=8, ymax=18,
            xtick distance=10,
            ytick distance=2,
            tick label style={font=\small},
            label style={font=\small},
            title style={font=\normalsize\bfseries},
        ]
        \addplot[color=magenta!70!black, line width=0.8pt]
            table[col sep=comma, x=tiempo, y=voltaje]{CSV/onda completa/oc_C2_220uF.csv};
        \end{axis}
    \end{tikzpicture}
    \caption{Tensi\'on de salida del rectificador de onda completa con $C_2 = 220\,\mu$F.}
    \label{fig:oc_c2_220uF}
\end{figure}

\begin{figure}[H]
    \centering
    \begin{tikzpicture}
        \begin{axis}[
            width=\textwidth,
            height=6cm,
            xlabel={Tiempo [ms]},
            ylabel={Tensi\'on [V]},
            title={Rectificador de onda completa --- $C_3 = 2200\,\mu$F},
            grid=both,
            grid style={line width=0.3pt, draw=gray!30},
            xmin=0, xmax=45,
            ymin=12, ymax=15,
            xtick distance=10,
            ytick distance=1,
            tick label style={font=\small},
            label style={font=\small},
            title style={font=\normalsize\bfseries},
        ]
        \addplot[color=magenta!70!black, line width=0.8pt]
            table[col sep=comma, x=tiempo, y=voltaje]{CSV/onda completa/oc_C3_2200uF.csv};
        \end{axis}
    \end{tikzpicture}
    \caption{Tensi\'on de salida del rectificador de onda completa con $C_3 = 2200\,\mu$F.}
    \label{fig:oc_c3_2200uF}
\end{figure}
